\chapter{Proposed Methods}
\label{sec:approach}

\section{System Overview}
The example compares a baseline retrieval configuration with a variant that
reranks an initial candidate set. Both configurations receive the same query
and search the same collection. The baseline returns its initial ordering,
whereas the variant applies an additional scoring function to the candidates.
Documents outside that candidate set cannot be promoted by the reranker, which
makes candidate selection an important part of the experimental description.

The implementation should record the candidate-set size, preprocessing rules,
and treatment of tied scores. These details can change the output even when
the high-level method name remains the same. Configuration files and software
versions should be saved with the result files. The appendix provides an
example checklist for this record, rather than prescribing a particular
software framework.

\section{Aggregation of Query Scores}
Let $s_i$ denote a selected evaluation score for query $i$, with $n$ queries
included in the comparison. The mean score is
\begin{equation}
  \label{eq:mean}
  \bar{s} = \frac{1}{n}\sum_{i=1}^{n}s_i.
\end{equation}
Equation~\eqref{eq:mean} assigns equal weight to each query. A different weighting
scheme would answer a different question and should be stated explicitly.
The metric used to obtain $s_i$ must also be defined; the aggregation formula
alone does not specify how relevance, rank or unjudged documents are handled.

The example retains both the mean and the individual scores. This allows the
analysis to distinguish a small improvement shared by many queries from a
large improvement concentrated in only a few cases. Neither pattern can be
reconstructed from the mean alone.

% Demonstration-only break.
\newpage
\section{Paired Comparison}
For each query, define the difference between the variant and baseline scores
using the same relevance judgements. Pairing the observations removes ambiguity
about which queries contributed to each system's mean. If a query is excluded,
the reason should be documented and the exclusion should be applied
consistently to both configurations.

A useful result file contains the query identifier, both scores and their
difference. Additional fields may describe the query group or judgement
coverage. These annotations support later analysis, but they should not be
chosen merely because they produce a favourable comparison. Decisions about
groups and exclusions belong in the experimental design wherever possible.

\subsection{Quality Checks}
Before aggregation, check for duplicate identifiers, missing outputs and
unexpected score ranges. The comparison should fail visibly when a required
system output is absent. Replacing an absent output with a score of zero may
be appropriate in some protocols, but it must be an explicit evaluation rule,
not an accidental consequence of a file-reading routine.

\section{Reproducibility}
The final report should distinguish deterministic settings from operations
that depend on a random seed. If repeated runs are used, explain what changes
between runs and how the resulting observations are summarised. Saving a seed
without recording the software and data versions is not a complete description
of the experimental environment.

The procedure should be checked on a small example before applying it to the
full collection. Manually inspecting a few query-level calculations can reveal
mismatched identifiers or inconsistent cutoffs. This check complements, rather
than replaces, automated validation of the complete result files.
